\documentclass[10pt,journal,compsoc]{IEEEtran}

% This manuscript is an original account of the implementation in this
% repository. It follows the IEEE journal structure without reproducing the
% wording of the source paper.
\usepackage{amsmath,amssymb,booktabs,graphicx,algorithm,algorithmic}
\usepackage{url}
\usepackage[hidelinks]{hyperref}

\newcommand{\FG}{\mathrm{FG}}
\newcommand{\AB}{\mathrm{AB}}

\title{Adaptive Candidate-Budget Beam Search for Folded-Subgraph Local Community Detection}

\author{Author Name\\
Department and Institution\\
City, Country\\
email@example.com}

\begin{document}
\maketitle

\begin{abstract}
Local community detection seeks the community containing a queried seed node
without requiring a global partition of the network. The folded-subgraph (FG)
method improves local modularity expansion by grouping tightly connected
neighbors before community growth. Its modularity stage, however, commits to a
single candidate at each iteration. A locally favorable candidate can
therefore prevent a later, more suitable expansion. This paper introduces a
validation-tuned adaptive candidate-budget beam-search extension of FG. The
extension preserves the folded-subgraph construction, density rule,
modularity equations, non-negative expansion constraint, and relaxation stage.
It changes only the search policy: several partial expansion paths are retained,
while a validated local budget limits the number of expensive folded candidates
examined at each state. On the generated LFR benchmark used by this
implementation, the method was selected on validation graphs with a beam
width of two, maximum depth five, learned path weight 0.10, and candidate
budget 64. Across all 3,000 eligible seeds in held-out LFR6--LFR9 graphs,
mean F1 increased from 0.621031 for the original implementation to 0.713990
for the adaptive method. Mean per-seed runtime was 180.392 ms for the adaptive
method and 1,323.625 ms for the original method in this Java implementation.
The improvement was not universal: LFR6 favored the original method, whereas
the adaptive method improved LFR7--LFR9. The results identify irreversible
greedy expansion as a substantive weakness while avoiding a claim of universal
dominance.
\end{abstract}

\begin{IEEEkeywords}
local community detection, folded subgraph, modularity expansion, beam
search, LFR benchmark, graph mining.
\end{IEEEkeywords}

\section{Introduction}
Community detection is commonly formulated as a global graph-partitioning
problem. In many applications, however, the user supplies one node and needs
only the local community surrounding that node. Local methods can reduce
memory and computation, but they must make decisions with incomplete context
and are sensitive to boundary edges.

The folded-subgraph method of Zhang and Bi groups neighboring nodes into dense
local structures before modularity expansion \cite{zhang2025fg}. This design
is attractive because it combines a local representation with a modularity
criterion and a relaxation step for noisy boundaries. The method nevertheless
uses a greedy modularity decision: after evaluating the available folded
subgraphs, it permanently selects the candidate with the greatest immediate
gain. That decision is not necessarily the best decision for the complete
expansion path.

This paper studies that failure mode and proposes a search-based extension.
The central idea is deliberately narrow. The original graph statistics and
folding rules are retained, but the modularity stage is allowed to preserve a
small number of alternatives. A learned local compatibility score is used only
as a secondary path-ranking term. To control the extra cost, candidates are
first ordered by inexpensive local connectivity information and only a
validation-tuned prefix is folded exactly.

The contributions are as follows:
\begin{enumerate}
  \item We identify irreversible candidate commitment as a distinct algorithmic
  weakness of the FG modularity stage.
  \item We implement a validation-tuned adaptive candidate-budget beam search
  that preserves the original folded-subgraph and relaxation mechanisms.
  \item We provide a Java implementation with deterministic preprocessing,
  restartable output, cached local computations, per-seed metrics, and figures.
  \item We evaluate every eligible seed in LFR6--LFR9, totaling 3,000 held-out
  seeds, and report both gains and regressions.
\end{enumerate}

The exact graph realizations used by the original paper are not distributed
with the paper supplied to this project. Accordingly, the experiments below
use deterministic generated LFR graphs with the paper's published macro
parameters. The results are implementation results on those realizations, not
claims that the paper's original Table III--V numbers have been reproduced.

\section{Background and Original FG Method}
\subsection{Local problem definition}
Let $G=(V,E)$ be an undirected graph and let $s\in V$ be a seed node. The goal
is to estimate the hidden community $C_s$ containing $s$ using local graph
access. The algorithm returns $\widehat{C}_s$. Evaluation compares the two
sets using recall, precision, and F1:
\begin{align}
 R &= \frac{|\widehat{C}_s\cap C_s|}{|C_s|},\\
 P &= \frac{|\widehat{C}_s\cap C_s|}{|\widehat{C}_s|},\\
 F_1 &= \frac{2RP}{R+P}.
\end{align}

\subsection{Folded-subgraph construction}
For a starting node $v$, the folding procedure initializes $F=\{v\}$ and
$R=N(v)$. At each step, it selects the candidate with the largest number of
links into $F$. Ties are resolved by smaller degree and then by node ID. The
candidate is tentatively added only when the induced density remains at least
$\gamma$:
\begin{equation}
 \rho(F)=\frac{2|E(F)|}{|F|(|F|-1)} \geq \gamma.
\end{equation}
The result is denoted by $F(v)$.

\subsection{Folding-community stage}
Starting from a current community $C$, the original method constructs folded
subgraphs for boundary candidates. It accepts a folded subgraph when at least
half of its nodes overlap the current community:
\begin{equation}
 |F(v)\cap C| \geq \frac{|F(v)|}{2}.
\end{equation}
The process repeats until a complete pass adds no node.

\subsection{Modularity stage}
For a community $C$, let $e_{in}$ denote internal edges and $e_{out}$ denote
edges leaving the community. The local modularity ratio is:
\begin{equation}
 M(C)=\frac{e_{in}}{e_{out}}.
\end{equation}
For a folded candidate $F$, the implementation follows the paper's local
quantities. Let $x$ count edges from new candidate nodes into $C$, let $y$
count the contribution of alternative nodes that connect into $C$, and let
$q$ count internal edges among new nodes. If $d$ is the sum of degrees of new
candidate nodes, the boundary adjustment is:
\begin{equation}
 z=d-x-\frac{q}{2}.
\end{equation}
The change used by the implementation is:
\begin{equation}
 \Delta M = \frac{e_{in}+x}{e_{out}+z-y}-M(C),
\end{equation}
when the denominator is positive. The original method selects the candidate
with maximum non-negative $\Delta M$ and discards the remaining alternatives.

\subsection{Relaxation stage}
The relaxation score combines first- and second-order contact with the current
community:
\begin{align}
 f_1(F,C)&=\frac{|N(F)\cap C|}{|C|},\\
 f_2(F,C)&=\frac{|NN(F)\cap C|}{|C|},\\
 f(F,C)&=f_1(F,C)+\beta f_2(F,C).
\end{align}
The implementation uses $\beta=0.5$ and accepts the best candidate when
$f(F,C)\geq h$. The full FG detector repeatedly applies folding-community
expansion, modularity expansion, and relaxation until no progress is made.

\section{Failure Analysis}
\subsection{Irreversible local choice}
The modularity stage ranks candidates by immediate gain. Suppose candidate
$a$ has a slightly larger $\Delta M$ than candidate $b$. The original method
chooses $a$ even if adding $b$ would expose a sequence of dense candidates that
eventually yields a better community. Because the current community changes
after each addition, the discarded candidate cannot be reconsidered.

\subsection{Capacity is not the primary issue}
The repository includes pointwise logistic, MLP, Bernoulli--Bernoulli RBM,
pairwise ranking, and shallow boosted-stump variants. The exploratory results
did not show a consistent gain from simply increasing predictor complexity.
This suggests that the main limitation is the irreversible search decision,
rather than the absence of a more expressive candidate classifier.

\subsection{Computational cost}
The original and unrestricted beam implementations repeatedly build folded
subgraphs and recompute local edge statistics. The final implementation uses
incremental folding connection counts, cached folded subgraphs, cached
community edge statistics, thread-local detector reuse, and resumable CSV
output. These optimizations do not alter the mathematical acceptance rules.

\section{Adaptive Beam-FG Method}
\subsection{Method overview}
The proposed method is called \emph{Validation-Tuned Adaptive Candidate-Budget
Beam-Search FG}, abbreviated AdaptiveBeam. It retains a list of partial
communities instead of a single modularity path. The state score is:
\begin{equation}
 S(C)=\sum_{t=1}^{d}\Delta M_t+\lambda p_t,
\end{equation}
where $p_t$ is a learned local compatibility probability and $\lambda$ is a
small secondary weight. The modularity gain remains the primary term.

\subsection{Candidate budget}
Before constructing an expensive folded subgraph, boundary candidates are
ordered by:
\begin{enumerate}
  \item decreasing number of links into the current community;
  \item increasing graph degree;
  \item increasing node identifier.
\end{enumerate}
Only the first $B$ candidates are evaluated exactly. The budget is not chosen
from held-out test data. It is selected on validation graphs from the set
$\{16,32,64,128\}$.

\subsection{Learned compatibility score}
The scorer is logistic regression over nine local features: folded density,
community overlap, first-order overlap, second-order overlap, modularity gain,
external edge ratio, folded size, mean degree, and internal candidate ratio.
Training uses LFR1--LFR5. Validation uses the first five eligible seeds in
LFR3--LFR5. The selected scorer contributes only $\lambda=0.10$ to the path
score; ground-truth labels are never used during held-out detection.

\begin{algorithm}[t]
\caption{Adaptive candidate-budget beam expansion}
\label{alg:adaptive}
\begin{algorithmic}[1]
\REQUIRE current community $C$, graph $G$, density $\gamma$, width $W$,
depth $D$, budget $B$
\STATE $\mathcal{B}\leftarrow\{(C,0)\}$
\STATE $C^*\leftarrow C$
\FOR{$d=1$ to $D$}
  \STATE $\mathcal{N}\leftarrow\emptyset$
  \FOR{each $(C_i,S_i)\in\mathcal{B}$}
    \STATE rank boundary candidates of $C_i$ by local priority
    \FOR{each of the first $B$ candidates $v$}
      \STATE construct $F(v)$ using the original folding rule
      \IF{$F(v)$ adds a node and $\Delta M(C_i,F(v))\geq0$}
        \STATE compute local features and $p(v)$
        \STATE add $(C_i\cup F(v),S_i+\Delta M+0.10p(v))$ to $\mathcal{N}$
      \ENDIF
    \ENDFOR
  \ENDFOR
  \STATE remove duplicate communities from $\mathcal{N}$
  \STATE retain the best $W$ states by score
  \IF{$\mathcal{N}=\emptyset$} \STATE \textbf{break} \ENDIF
  \STATE $\mathcal{B}\leftarrow\mathcal{N}$
  \STATE update $C^*$ with the highest-scoring state
\ENDFOR
\RETURN $C^*$
\end{algorithmic}
\end{algorithm}

The returned community is then passed through the unchanged relaxation stage.
Thus, the new method changes the search policy, not the definition of a valid
folded subgraph or the paper's modularity equations.

\section{Experimental Design}
\subsection{Datasets}
The implementation encodes the LFR parameter rows reported in the supplied
paper. LFR1--LFR5 contain 200 nodes with average degree 7, maximum degree 15,
and mixing coefficients from 0.1 to 0.5. LFR6 contains 1,000 nodes with
average degree 15, maximum degree 50, and maximum community size 50. LFR7
uses the same graph scale with maximum community size 300. LFR8 and LFR9
contain 500 nodes with average degree 7 and maximum degree 15. All rows use
the specified degree and community exponents and non-overlapping communities.

The supplied paper does not provide the complete original graph files and
random seeds. The repository therefore uses deterministic generated
realizations with the published macro parameters. This distinction is
essential when comparing the present numbers with the paper's tables.

\subsection{Protocol}
The model is trained before held-out testing. Configuration search is limited
to validation data. Once the candidate budget and path parameters are frozen,
LFR6--LFR9 remain untouched until final evaluation. Every eligible node in the
held-out graphs is used as a seed. The final evaluation contains 3,000 seeds
per method.

\subsection{Metrics}
For each seed, recall, precision, F1, and runtime are recorded. The summary
reports arithmetic means over seeds in each dataset. Paired comparisons use
the same seed for the two methods. The report does not interpret a positive
mean difference as proof of statistical significance without a separately
specified confidence procedure.

\section{Results}
\subsection{Validation}
The validation sweep selected $B=64$:
\begin{table}[t]
\caption{Candidate-budget validation results.}
\label{tab:validation}
\centering
\begin{tabular}{rr}
\toprule
Budget & Mean validation F1\\
\midrule
16 & 0.339110\\
32 & 0.337700\\
64 & \textbf{0.369291}\\
128 & 0.368839\\
\bottomrule
\end{tabular}
\end{table}

\subsection{All-seed comparison}
Table~\ref{tab:f1} reports the complete held-out comparison.
\begin{table}[t]
\caption{Mean F1 over all eligible held-out seeds.}
\label{tab:f1}
\centering
\begin{tabular}{lrrrr}
\toprule
Dataset & Seeds & FG & Full beam & AdaptiveBeam\\
\midrule
LFR6 & 1000 & 0.754533 & 0.686188 & 0.686188\\
LFR7 & 1000 & 0.338073 & 0.586587 & \textbf{0.593505}\\
LFR8 & 500 & 0.737770 & 0.851961 & \textbf{0.851974}\\
LFR9 & 500 & 0.803206 & 0.872581 & 0.872581\\
All & 3000 & 0.621031 & 0.711682 & \textbf{0.713990}\\
\bottomrule
\end{tabular}
\end{table}

\begin{table}[t]
\caption{Overall recall, precision, F1, and runtime.}
\label{tab:overall}
\centering
\begin{tabular}{lrrrr}
\toprule
Method & Recall & Precision & F1 & ms/seed\\
\midrule
FG & 0.948242 & 0.538054 & 0.621031 & 1323.625\\
Full beam & 0.889162 & 0.647175 & 0.711682 & 249.489\\
AdaptiveBeam & 0.892300 & 0.648429 & \textbf{0.713990} & \textbf{180.392}\\
\bottomrule
\end{tabular}
\end{table}

AdaptiveBeam improves mean F1 over FG by 0.092959, or approximately 14.97
percent relative to the FG mean. It improves over unrestricted beam search by
0.002308 while reducing mean runtime by approximately 27.7 percent. In the
paired seed comparison with FG, AdaptiveBeam wins 1,253 seeds, ties 1,235,
and loses 512.

\begin{figure}[t]
\centering
\includegraphics[width=\columnwidth]{../output/adaptive-beam-results-full/figures/adaptive-beam-vs-fg-f1.png}
\caption{F1 comparison over LFR6--LFR9. The adaptive method preserves the
beam improvement on LFR7--LFR9 while remaining close to unrestricted beam.}
\label{fig:f1}
\end{figure}

\section{Discussion}
The most important result is not only the aggregate mean. The dataset-wise
pattern reveals the operating regime of the new method. LFR7 is the clearest
case: the original greedy method obtains F1 0.338073, while unrestricted and
adaptive beam obtain 0.586587 and 0.593505. This supports the hypothesis that
greedy commitment is harmful when the local graph contains several plausible
expansion routes. LFR8 and LFR9 show the same direction with smaller
differences.

LFR6 is a counterexample. FG obtains F1 0.754533, while both beam variants
obtain 0.686188. A beam search can preserve an attractive alternative that is
not favorable for this graph regime, and the learned secondary score cannot
guarantee correction. The method should therefore be presented as a robust
search extension with regime-dependent behavior, not as a universally superior
replacement.

The precision and recall averages also explain the F1 change. FG has higher
recall overall but much lower precision. Beam methods reduce recall somewhat
while improving precision substantially. The resulting balance produces the
higher F1 values. This is consistent with the paper's motivation that local
community methods must control boundary noise rather than simply absorb more
neighbors.

\section{Threats to Validity and Limitations}
First, the original graph realizations and random seeds used by the paper are
not available in this repository. The LFR results reproduce the published
parameter regimes, not the authors' exact realizations. Second, the complete
experiment is a benchmark study on four held-out LFR settings; it is not a
substitute for evaluation on the paper's real-world networks. Third, the
candidate budget was selected from a small validation sweep. A larger study
could investigate budget stability across graph families. Fourth, runtime
depends on the Java runtime, hardware, graph representation, and cache state.
Finally, the reported paired counts are descriptive; a formal confidence
interval or paired randomization test should be added before a publication
claim about statistical significance.

\section{Reproducibility}
The source code, generated data, configuration, per-seed CSV, summary, and
figure are stored in the public repository:
\begin{center}
\url{https://github.com/sankar1912/fg-based-local-community-detection-java}
\end{center}

The main commands are:
\begin{verbatim}
javac -Xlint:all -d out *.java
java -cp out foldedsubgraph.RigorousEvaluationTest
java -cp out foldedsubgraph.AdaptiveBeamExperiment \\
  datasets/generated-lfr output/adaptive-beam-results-full 0
\end{verbatim}

The final configuration is recorded in
\texttt{output/adaptive-beam-results-full/experiment-config.txt}. The complete
summary is in \texttt{summary.csv}, and the 9,000 per-seed rows are in
\texttt{per-seed.csv}. The implementation also records validation results so
that the selected budget can be audited.

\section{Conclusion}
This paper presented AdaptiveBeam, a validation-tuned adaptive
candidate-budget beam-search extension of folded-subgraph local community
detection. The method targets a specific weakness of the original FG
algorithm: its irreversible selection of one immediate modularity winner. By
retaining a small number of expansion paths and restricting expensive folding
to a validated candidate budget, the method improves the overall F1 score and
reduces runtime relative to unrestricted beam search in the complete
evaluation.

Across 3,000 held-out seeds, AdaptiveBeam reaches mean F1 0.713990 compared
with 0.621031 for FG. The result is encouraging but conditional: LFR6 remains
a regression, and the graph realizations are generated from published LFR
parameters rather than recovered from the authors' private experiment files.
The appropriate conclusion is that adaptive multi-path expansion is a
promising improvement and that irreversible greedy commitment is a measurable
weakness, not that the new method is universally optimal.

\begin{thebibliography}{00}
\bibitem{zhang2025fg}
M. Zhang and W. Bi, ``A local community detection method based on folded
subgraph,'' \emph{IEEE Transactions on Knowledge and Data Engineering},
vol. 37, no. 7, pp. 3869--3880, 2025.

\bibitem{lfr2008}
A. Lancichinetti, S. Fortunato, and F. Radicchi, ``Benchmark graphs for
testing community detection algorithms,'' \emph{Physical Review E}, vol. 78,
no. 4, Art. no. 046110, 2008.

\bibitem{newman2006}
M. E. J. Newman, ``Modularity and community structure in networks,''
\emph{Proceedings of the National Academy of Sciences}, vol. 103, no. 23,
pp. 8577--8582, 2006.

\bibitem{freund1997}
Y. Freund and R. E. Schapire, ``A decision-theoretic generalization of
on-line learning and an application to boosting,'' \emph{Journal of Computer
and System Sciences}, vol. 55, no. 1, pp. 119--139, 1997.
\end{thebibliography}

\end{document}
